%=============================================================================!
% 7.E  EXERCISES                                                              !
%=============================================================================!
\unnumbered{Exercises}
\label{sec:ha-exercises}

The exercises run from questions of understanding, through calculations,
to short derivations marked `show that'. Full solutions are in
\hyperref[sec:ha-solutions]{`Solutions to the exercises'} at the end of
the book, and Notebook~07 checks every numerical answer.

\begin{exercise}[a transform with a logarithm]
\label{ex:ha-exp}
Show that the Legendre transform of $f(v) = \mathrm{e}^v$ is $f^*(p) =
p\ln p - p$, for $p > 0$, and that its slope is $v$.
\end{exercise}

\begin{exercise}[two rates in one term]
\label{ex:ha-cross}
A kinetic energy $K = \frac12(a\dot{q}_1^2 + 2b\,\dot{q}_1\dot{q}_2 +
c\,\dot{q}_2^2)$, with constants $a$, $b$ and $c$, contains a product of
two rates. Show that $p_1\dot{q}_1 + p_2\dot{q}_2 = 2K$, so that again
$H = K + U$.
\end{exercise}

\begin{exercise}[a hanging spring]
\label{ex:ha-vertical}
A block of mass $m$ hangs from a spring of stiffness $k$, with $q$ its
distance below the point where the spring has its natural length, so that
$U = \frac12 kq^2 - mgq$. Write $H$ and Hamilton's equations, and find the
state that does not change.
\end{exercise}

\begin{exercise}[a ball in a vertical plane]
\label{ex:ha-plane}
For $\Lag = \frac12 m(\dot{x}^2 + \dot{y}^2) - mgy$, find $p_x$, $p_y$
and $H$, write Hamilton's equations, and say which momentum is conserved
and why.
\end{exercise}

\begin{exercise}[the puck]
\label{ex:ha-puck}
Show that the puck of \cref{sec:la-polar}, with $\Lag = \frac12 m(\dot{r}^2
+ r^2\dot\theta^2) - U(r)$, has $H = p_r^2/(2m) + p_\theta^2/(2mr^2) +
U(r)$, and that Hamilton's equations give the conservation of $p_\theta$
and the motion of $r$ in the effective potential \cref{eq:la-effective}.
\end{exercise}

\begin{exercise}[the bead from its Hamiltonian]
\label{ex:ha-bead}
Show that Hamilton's equations for $H = p^2/\bigl(2m(1 + h'^2)\bigr) +
mgh(x)$ give the bead's equation \cref{eq:la-bead}.
\end{exercise}

\begin{exercise}[over the top]
\label{ex:ha-separatrix}
For the pendulum of \cref{fig:ha-portrait}, find the momentum and the
speed of the bob at the bottom of the separatrix.
\end{exercise}

\begin{exercise}[the last ten degrees]
\label{ex:ha-creep}
Using $\dot\varphi \approx -\omega_0\varphi$, how long does the bob of
\cref{fig:ha-portrait}, on its separatrix, take to go from \ang{10} short
of the top to \ang{1} short?
\end{exercise}

\begin{exercise}[a product rule for brackets]
\label{ex:ha-product}
Show that $\pb{A}{BC} = \pb{A}{B}\,C + B\,\pb{A}{C}$.
\end{exercise}

\begin{exercise}[angular momentum in Cartesian coordinates]
\label{ex:ha-lz}
A body moves in a plane with $H = (p_x^2 + p_y^2)/(2m) + U(r)$, where $r
= \sqrt{x^2 + y^2}$. Show that $L_z = xp_y - yp_x$ has $\pb{L_z}{H} = 0$.
\end{exercise}

\begin{exercise}[the momentum of the spring]
\label{ex:ha-spring-p}
Use the series \cref{eq:ha-series} to find $p(t)$ for the block on a
spring.
\end{exercise}

\begin{exercise}[a shear]
\label{ex:ha-shear}
A body of \SI{1}{\kilogram} flies freely. Where does the flow over
\SI{2}{\second} carry the four corners of the square of starting states
$0 \le q_0 \le \SI{0.1}{\metre}$, $0 \le p_0 \le
\SI{0.1}{\kilogram\metre\per\second}$, and what is the area of the
figure they make?
\end{exercise}

\begin{exercise}[half the area]
\label{ex:ha-halving}
With drag at $\gamma = \SI{0.5}{\per\second}$, how long does a blob of
states take to lose half its area?
\end{exercise}

\begin{exercise}[two carts]
\label{ex:ha-carts}
For the two carts of \cref{ex:la-carts}, $\Lag = \frac12 m_1\dot{x}_1^2 +
\frac12 m_2\dot{x}_2^2 - \frac12 k(x_2 - x_1 - \ell)^2$, use Noether's
theorem to show that $m_1\dot{x}_1 + m_2\dot{x}_2$ is conserved.
\end{exercise}

\begin{exercise}[under gravity]
\label{ex:ha-gravity}
A group of atoms with an energy $U_{\mathrm{int}}$ that is unchanged by
moving or turning the whole group falls in a uniform field of gravity
along $-z$, so $U = U_{\mathrm{int}} + \sum_im_igz_i$. Which components
of the total momentum, and of the angular momentum about the origin, are
conserved?
\end{exercise}

\begin{exercise}[the torque about any point]
\label{ex:ha-torque-point}
Show that when the forces on a group add to zero, their torque is the
same about every point, and find the torque of the two atoms of
\cref{fig:ha-noether}(b) about the centre of the box, $(5,
5)$~\si{\angstrom}.
\end{exercise}

\begin{exercise}[which can be reversed?]
\label{ex:ha-even}
For which of $H_1 = p^2/(2m) + U(q)$, $H_2 = (p - c)^2/(2m)$ with a
constant $c \ne 0$, and $H_3 = p^2/(2m) + \frac12 kq^2 + bqp$ with a
constant $b \ne 0$ does the argument of \cref{sec:ha-reversal} show that
reversing the momentum retraces the path?
\end{exercise}

\begin{exercise}[doubling the energy]
\label{ex:ha-doubling}
The forward Euler method follows the block of \cref{fig:ha-maps} with
$\dt = \SI{0.005}{\second}$. After how many steps, and how long, has the
energy doubled?
\end{exercise}

\begin{exercise}[the other order]
\label{ex:ha-other}
The other symplectic Euler method moves the position first, $q' = q +
\dt\,p/m$, then the momentum with the force at the new position, $p' =
p - \dt\,kq'$. Show that for the spring its Jacobian determinant is 1.
\end{exercise}

\begin{exercise}[the energy the method keeps]
\label{ex:ha-shadow}
Show that a step of the symplectic Euler method of
\cref{sec:ha-symplectic} leaves $\tilde{H} = p^2/(2m) + \frac12 kq^2 -
\bigl(k\dt/(2m)\bigr)\,qp$ unchanged for the spring.
\end{exercise}
